\thispagestyle{empty}
\selectlanguage{hebrew}

\begin{center}
  {\Large\bfseries \DocumentTitleHebrew\par}
  \vspace{0.4cm}
  {\large \AuthorNameHebrew\par}
  \vspace{0.4cm}
  {\large עבודת גמר לתואר \DegreeNameHebrew\par}
  \vspace{0.2cm}
  {\large אוניברסיטת בן-גוריון בנגב\par}
  \vspace{0.2cm}
  \selectlanguage{english}{\large \ThesisYear\par}
  \selectlanguage{hebrew}
  \vspace{0.9cm}
  {\Large\bfseries תקציר\par}
\end{center}

אנו חוקרים הסקה בייסיאנית מוכללת מעל מדיניויות דטרמיניסטיות מלאות בתהליך
החלטה מרקובי סטוכסטי בעל אופק סופי. התפלגות המטרה נקבעת לפי תוחלת התגמול:
תחילה ממצעים על האקראיות בסביבה, ורק לאחר מכן מפעילים את הפונקציה המעריכית.
הגדרה זו, המבוססת על עבודות קודמות, שונה ממתן משקל מעריכי לתגמול שהתקבל
במסלול אקראי יחיד.

השיטה המעשית הנבחנת לומדת התפלגות מפורקת מעל טבלאות מדיניות מלאות.
אלגוריתם \foreignlanguage{english}{PPO} מספק אתחול בעל עלות מדווחת בנפרד.
לאחריו מתבצע עדכון וריאציוני ישיר בעזרת תגמולים מהסימולטור, הפחתת שונות
ואיבר אנטרופיה המחושב בדיוק. אנו מצדיקים את אומד הגרדיאנט ומשווים את
תוחלתו לחישוב מדויק בבעיות קטנות. מסלול אחד לכל מדיניות נדגמת מספיק
לתקינות האומד, אך חלוקת התקציב בין מדיניויות שונות לבין מסלולים חוזרים
משפיעה על הלמידה. בכל עדכון נדגמות מדיניויות חדשות; בזמן ביצוע נשמרת
המדיניות שנבחרה לכל אורך האפיזודה. אין בשיטה זו שלב של שכפול חלקיקים.

התכנון הניסויי מבחין בין תקינות הגרדיאנט, שגיאת משפחת הקירוב וביצועי
הבקרה. תקציב האינטראקציות כולל גם את האתחול וגם את העדכון הווריאציוני.
ההערכה מפרידה בין התחייבות למדיניות מלאה לבין הגרלת פעולה מחדש בכל ביקור.
ניסויי פיתוח ראשוניים בסביבות רשת מצביעים על תועלת לדגימת יותר מדיניויות
במקום חזרה מרובה על כל מדיניות, ועל ביצועים קרובים לבקרת
\foreignlanguage{english}{PPO} מסייעת בחקירה בניסוי בינוני אחד.
בניסוי הקשה, הקצאה קטנה מדי לשלב החקירה לא יצרה מסלולים מוצלחים.
חלוקה מחדש של אותו תקציב שיפרה את הביצוע, אך לא ביטלה את הפער בין כללי
הביצוע. אלו תוצאות פיתוח, ולא הוכחה לעליונות כללית.

העבודה כוללת גם שיטות דגימה חלופיות, בדיקות תקינות ותיעוד של ניסיונות
שלא הצליחו. השוואה מלאה מול שיטות הבסיס, חזרות בלתי תלויות, מפות שלא
שימשו לפיתוח וניתוח שגיאת הקירוב עדיין נדרשים להשלמת המסקנות.

\selectlanguage{english}
